\BourbakiExercises{Exercises}

\begin{enumerate}
\def\labelenumi{\arabic{enumi})}
\item
  Let A be a semisimple ring, and let M be a finitely generated A-module. For any submodule N of M, we denote the orthogonal of N in \(M^{*}\) by \(N'\) . Prove that the orthogonal of \(N'\) in M is equal to N and that we have \(\text{long}_{\text{A}}(\text{N}) + \text{long}_{\text{A}}(\text{N}') = \text{long}_{\text{A}}(\text{M})\) . The mapping \(N \mapsto N'\) is a bijection from the set of submodules of M to the set of submodules of \(M^{*}\) . If L is another submodule of M, then the orthogonal of \(L \cap N\) is \(L' + N'\) .
\item
  Let A be a ring. Prove that the \((A,A)\) -bimodule \(_{s}A_{d}\) is semisimple if and only if A is isomorphic to a finite product of quasi-simple rings.
\item
  We say that a (left or right) ideal \(\mathfrak{a}\) of a ring is nilpotent if there exists an integer \(n\) such that every product of \(n\) elements of \(\mathfrak{a}\) is zero.
\end{enumerate}

\begin{enumerate}
\def\labelenumi{\alph{enumi})}
\item
  Let \(\mathfrak{a}\) be a nilpotent left ideal of A. Prove that the two-sided ideal \(\mathfrak{aA}\) generated by \(\mathfrak{a}\) is nilpotent.
\item
  We say that a ring A is semiprime if it does not contain any nonzero nilpotent two-sided ideals. We say that a two-sided ideal \(\mathfrak{a}\) of A is semiprime if the quotient ring A/ \(\mathfrak{a}\) is semiprime. Prove that the intersection of a family of semiprime (two-sided) ideals is a semiprime ideal and that a semiprime ideal contains every nilpotent ideal of A.
\item
  A prime (two-sided) ideal (VIII, p.~133, Exercise 15) is semiprime. A two-sided ideal a is semiprime if and only if it is the intersection of the prime ideals that contain it. (If \(a \notin a\) , consider an ideal that is maximal among the two-sided ideals containing a and does not contain any powers of a.)
\end{enumerate}

¶ 4) Let A be a semiprime left Noetherian ring (Exercise 3) and S the set of its cancelable elements (cf.~VIII, p.~75, Exercise 8).

\begin{enumerate}
\def\labelenumi{\alph{enumi})}
\item
  Prove that every nonzero (left or right) ideal of A contains an element that is not nilpotent. (Let a be a right ideal of A consisting of nilpotent elements. Let x be a nonzero element of a whose left annihilator is maximal. Observe that for every element a of A and every integer \(n \geqslant 1\) such that \((xa)^{n} \neq 0\) , the left annihilator of \((xa)^{n}\) is equal to that of x, and deduce that we have xax = 0.)
\item
  Prove that S allows a calculus of left fractions on A (VIII, p.~18, Exercise 18). (For \(a \in \mathrm{A}\) and \(s \in \mathrm{S}\) , prove that the set of elements \(x\) of A such that \(xa \in \mathrm{As}\) is an irreducible left ideal of A (VIII, p.~75, Exercise 7), and conclude using \(a\) ) and Exercise 8, \(d\) ) of VIII, p.~75.)
\end{enumerate}

\(c\) ) Prove that the ring \(\mathrm{S}^{-1}\mathrm{A}\) is semisimple (``Goldie's theorem'' (1)). (Let \(\mathfrak{b}\) be an irreducible left ideal of \(\mathrm{S}^{-1}\mathrm{A}\) . Show that \(\mathrm{A} \cap \mathfrak{b}\) is an irreducible left ideal of A, and then use Exercise 8, d) of VIII, p.~75 to show that \(\mathfrak{b} = \mathrm{S}^{-1}\mathrm{A}\) . Conclude using Exercise 7, b) of VIII, p.~75.)

\( ^{(1)} \) For other approaches to this result, see X, §1, p.~171, exercice 25 and Comm. Alg., II, §2, p.~135, Exercise 26.

\begin{enumerate}
\def\labelenumi{\arabic{enumi})}
\setcounter{enumi}{4}
\tightlist
\item
  \begin{enumerate}
  \def\labelenumii{\alph{enumii})}
  \tightlist
  \item
    Let A be a semiprime ring such that every set of monogenous left ideals admits a minimal element. Prove that A is semisimple. (Using Exercise 7 of VIII, p.~130, construct sequences \((e_n)\) and \((f_n)\) of idempotents in A such that the ideals \(Ae_n\) are minimal and that we have \(A f_{n-1} = Ae_n \oplus Af_n\) , and show that \(f_n\) is zero for \(n\) sufficiently large.)
  \end{enumerate}
\end{enumerate}

\begin{enumerate}
\def\labelenumi{\alph{enumi})}
\setcounter{enumi}{1}
\tightlist
\item
  Let K be a commutative field, V an infinite-dimensional right vector space over K, and A the K-subalgebra of \(\operatorname{End}_{\mathrm{K}}(\mathrm{V})\) generated by \(1_{V}\) and the endomorphisms of finite rank. Prove that every left ideal of A contains a minimal ideal.
\end{enumerate}

\begin{enumerate}
\def\labelenumi{\arabic{enumi})}
\setcounter{enumi}{5}
\tightlist
\item
  Let A be a semisimple commutative algebra of finite rank over a commutative field K, and let \(K_{i}\) , for \(i \in I\) , be its simple factors, which are finite extensions of K. Let B be a subalgebra of A and \(E_{j}\) ( \(j \in J\) ) its simple factors. Prove that there exist pairwise disjoint nonempty subsets \((\Lambda_{j})_{j \in J}\) of I such that \(E_{j}\) injects into \(\prod_{i \in \Lambda_{j}} K_{i}\) for \(j \in J\) . For every \(i \in \Lambda_{j}\) , \(pr_{i}\) induces an isomorphism from \(E_{j}\) to a subfield \(E_{ij}\) of \(K_{i}\) . Study the converse. Deduce that if each of the \(K_{i}\) is separable over K, then A contains only finitely many subalgebras (cf.~V, §6, No.~4, p.~30, Proposition 3).
\end{enumerate}

